\section{Finite coefficient formulas and explicit mixed families}
\label{sec:coefficients}

\begin{corollary}[All mixed $u_3$ coefficients]\label{cor:coefficients}
Let $a,k,b\ge0$, $n\ge2$, and $m=a+2k+3b+4n$. If $k+n$ is odd,
then $[u_0^au_1^ku_2^bu_3^n]N_m=0$. If $k+n$ is even, then
\begin{equation}\label{eq:coefficient-integral}
\begin{split}
 &[u_0^au_1^ku_2^bu_3^n]N_m\\
 &=-\frac{m!}{9216\,a!2^kk!b!n!24^b48^{n-2}}
 \int_0^1v^{k+n-2}(1-v^2)^3
       (1+3v^2)^b(1+v^2)^{n-2}\dd v.
\end{split}
\end{equation}
Equivalently, the finite rational formula \eqref{eq:preview} holds.
Every coefficient not forced to vanish by parity is strictly negative.
\end{corollary}
\begin{proof}
Expand $e^{u_0t}$, $e^{yv}$, $e^{zh(v)}$, and $\Erem(wk(v))$ in
\eqref{eq:main}. The conversion from $y^k$ to $u_1^k$ contributes
$2^{-k}$, and conversion to $N_m$ contributes $m!$. The remaining
integrand is odd precisely when $k+n$ is odd; otherwise its average on
$[-1,1]$ is its integral on $[0,1]$. This proves
\eqref{eq:coefficient-integral}. Expand the two binomials and use
\begin{equation}\label{eq:beta-four}
 \int_0^1v^K(1-v^2)^3\dd v
 =\frac{48}{(K+1)(K+3)(K+5)(K+7)},\qquad K>-1.
\end{equation}
Formula~\eqref{eq:beta-four} follows by $x=v^2$ and three elementary
integrations of the beta integral, or by expanding $(1-v^2)^3$.
Here $K=k+n-2+2j+2\ell\ge0$, so every denominator is positive.
This proves the finite formula and the sign assertion.
\end{proof}

\subsection{The source square-sector question}
For $n=2$ and even $k$, \eqref{eq:preview} becomes
\begin{equation}\label{eq:square-answer}
\boxed{\begin{aligned}
 &[u_0^au_1^ku_2^bu_3^2]N_{a+2k+3b+8}\\
 &=-\frac{(a+2k+3b+8)!}{384\,a!2^kk!b!24^b}
 \sum_{j=0}^b
 \frac{\binom bj3^j}
 {(k+2j+1)(k+2j+3)(k+2j+5)(k+2j+7)}.
\end{aligned}}
\end{equation}
It is zero for odd $k$. Thus every $b>0$ in the specific question
\eqref{eq:source-question} is covered; no weight-dependent linear
system remains to be solved.

Define the normalized square coefficient for even $k$ by
\[
 g_b(k)=k![y^kz^bw^2]F(y,z,w).
\]
Then
\begin{equation}\label{eq:gb-density}
 g_b(k)=-\frac1{18432b!24^b}
 \int_0^1v^k(1-v^2)^3(1+3v^2)^b\dd v.
\end{equation}
For example,
\begin{align}
 g_0(k)&=-\frac1{384(k+1)(k+3)(k+5)(k+7)},\label{eq:g0}\\
 g_1(k)&=-\frac1{2304(k+1)(k+5)(k+7)(k+9)},\label{eq:g1}\\
 g_2(k)&=-\frac{(k+2)(k+6)}{27648\prod_{j=0}^5(k+2j+1)},\label{eq:g2}\\
 g_3(k)&=-\frac{k^3+15k^2+62k+63}{497664\prod_{j=0}^6(k+2j+1)}.
 \label{eq:g3}
\end{align}
These rational functions are asserted as coefficient formulas at even
nonnegative integers $k$. Their analytic continuation in $k$ is also
provided by the beta integral for $\Re k>-1$, but the coefficient at an
odd power of $y$ remains zero.

\subsection{Nonlinear sectors beyond the square}
The complete formula also gives genuinely coupled cubic and quartic
sectors. A small sample, with no $u_0$ factor, is
\begin{center}
\begin{tabular}{rll}
\toprule
Weight & Functional monomial & Coefficient in $N_m$\\
\midrule
8 & $u_3^2$ & $-1$\\
11 & $u_2u_3^2$ & $-55$\\
12 & $u_1^2u_3^2$ & $-165$\\
14 & $u_2^2u_3^2$ & $-3640$\\
14 & $u_1u_3^3$ & $-3185/3$\\
16 & $u_3^4$ & $-10325/3$\\
17 & $u_1u_2u_3^3$ & $-1023400/3$\\
19 & $u_2u_3^4$ & $-1639225$\\
\bottomrule
\end{tabular}
\end{center}
All coefficients with an additional power $u_0^a$ follow immediately
from the Appell factorial ratio. The rational denominators in this
basis need not be integral, even though the Bell polynomial before
normal reduction has integral coefficients.

The formulas also reveal a useful structural distinction. The pure
highest-index-two sector has positive beta density, whereas the entire
highest-index-three sector has negative beta density, subject to parity.
This observation motivates the all-order conjecture in
Section~\ref{sec:conjecture}; it does not by itself establish that
conjecture.
